%
% --- inline annotations
%
\ifdefined\XeTeXversion\else\pdfoutput=1\fi % only force PDF output under pdfTeX; breaks hyperref's driver detection under XeTeX
\usepackage[dvipsnames, table]{xcolor}
\usepackage{color}
% \usepackage[nointegrals]{wasysym} 
\usepackage{tikz,pgfplots,pgfplotstable}
% \usepackage{stix}
\usepackage{float}
\usepackage{amsthm}
\usepackage{diagbox}
\usepackage{xspace}
\usepackage{multirow}
\usepackage{tabularx}
% xr-hyper is the hyperref-aware variant of xr; load before hyperref
\usepackage{xr-hyper}
\usepackage{caption}
\usepackage{pifont}
\usepackage{amsmath}
\makeatletter\@ifpackageloaded{hyperref}{}{\usepackage{hyperref}}\makeatother
\usepackage{wrapfig}
\usepackage{enumitem}

% \newcommand{\vmduc}[1]{\textcolor{red}{}}
\newcommand{\vmduc}[1]{\textcolor{red}{\small VMDuc: #1}}
% \newcommand{\hh}[1]{\textcolor{blue}{#1}}
\newcommand{\hh}[1]{#1}
\newcommand{\E}[1]{\mathbb{E}_{P_t}\left[#1\right]}
\newcommand{\Var}[1]{\mathrm{Var}_{P_t}\left(#1\right)}

% Define 
\newcommand{\method}[0]{PeTTA\xspace}
\newcommand{\setting}[0]{recurring\xspace}
\newcommand{\Setting}[0]{Recurring\xspace}


\newcommand{\cmark}{\ding{51}}%
\newcommand{\xmark}{\ding{55}}%
\definecolor{ClrHighlight}{RGB}{210, 255, 255}
\definecolor{green01270}{RGB}{80, 200, 120}
\definecolor{ForestGreen}{RGB}{34,139,34}
\definecolor{darkgray}{RGB}{113, 121, 126}

\newcommand{\red}[1]{{\color{red}#1}}
% \newcommand{\todo}[1]{{\color{red}#1}}
% \newcommand{\TODO}[1]{\textbf{\color{red}[TODO: #1]}}

\newtheorem{thm}{Theorem}[subsection]
\newtheorem{theorem}{Theorem}
\newtheorem{lemma}{Lemma}
\newtheorem{corollary}{Corollary}
\newtheorem{assumption}{Assumption}
\newtheorem{definition}{Definition}

\newenvironment{shortthm}{\textbf{Theorem.}\space\em}{\par}

% --- disable by uncommenting  
% \renewcommand{\TODO}[1]{}
% \renewcommand{\todo}[1]{#1}
\newcommand{\changes}[1]{\textcolor{black}{#1}}
\newcommand{\justification}[1]{\textcolor{blue}{#1}}
\newcommand{\remove}[1]{\textcolor{red}{#1}}
% Setup embedded hyperlinks
\definecolor{cvprblue}{rgb}{0,0.20,0.74}
\hypersetup{
    colorlinks=true,
    linkcolor=cvprblue,
    filecolor=magenta,      
    urlcolor=cyan,
    % pdftitle={Overleaf Example},
    % pdfpagemode=FullScreen,
    }
    
\newcommand{\Sec}[1]{{Sec.}~#1}
\newcommand{\Tab}[1]{{Tab.}~#1}
\newcommand{\Fig}[1]{{Fig.}~#1}
\newcommand{\Appdx}[1]{{Appdx.}~#1}
\newcommand{\Eq}[1]{{Eq.}~#1}
\newcommand{\Alg}[1]{{Alg.}~#1}

\usepackage{algorithm} 
\usepackage{algorithmic}  
\usepackage[algo2e]{algorithm2e} 
\usetikzlibrary{fit,calc}
% For highlighting a segment of code block
%define a marking command
\newcommand*{\tikzmk}[1]{\tikz[remember picture,overlay,] \node (#1) {};\ignorespaces}
%define a boxing command, argument = colour of box
\newcommand{\boxit}[1]{\tikz[remember picture,overlay]{\node[yshift=3pt,fill=#1,opacity=.25,fit={(StartHere)($(EndHere)+(.95\linewidth,.8\baselineskip)$)}] {};}\ignorespaces}
%define some colours according to algorithm parts (or any other method you like)
\colorlet{codehighlightpink}{red!40}
\colorlet{codehighlightblue}{cyan!60}

\newcommand{\dataname}{SafeGrad\xspace}

\definecolor{bestresult}{RGB}{198,224,180}  % Light green for best
\definecolor{worstresult}{RGB}{248,203,173} % Light orange for worst
\definecolor{highlight}{RGB}{255,242,204}    % Light yellow for notable

% \usepackage{floatrow}

\usepackage[utf8]{inputenc} % allow utf-8 input
\usepackage[T1]{fontenc}    % use 8-bit T1 fonts
\makeatletter\@ifpackageloaded{hyperref}{}{\usepackage{hyperref}}\makeatother       % hyperlinks
\usepackage{url}            % simple URL typesetting
\usepackage{booktabs}       % professional-quality tables
\usepackage{amsfonts}       % blackboard math symbols
\usepackage{nicefrac}       % compact symbols for 1/2, etc.
\usepackage{microtype}      % microtypography
\usepackage{xcolor}         % colors
\usepackage{xspace}

\usepackage{tikz}             % The engine for the diagram
\usepackage{xcolor}           % For defining the safety level colors (Green/Red/etc.)
\usepackage{caption}          % For professional caption styling
\usepackage{subcaption}       % If you plan to have Figure 1a and 1b
\usepackage{amsmath, amssymb} % For any math symbols in the rules
\usetikzlibrary{shapes.geometric} % For the Category Wheel and rounded boxes
\usetikzlibrary{arrows.meta}      % For professional, modern arrowheads
\usetikzlibrary{positioning}      % For relative placement (e.g., "right=of category")
\usetikzlibrary{calc}             % To calculate coordinates for the 'Safety Gap' arrow
\usetikzlibrary{shadows}          % Optional: to give boxes a subtle drop-shadow
\definecolor{safe_green}{RGB}{220, 245, 220}
\definecolor{low_yellow}{RGB}{255, 250, 205}
\definecolor{mid_orange}{RGB}{255, 225, 180}
\definecolor{high_red}{RGB}{255, 210, 210}
\definecolor{safety_gap_red}{RGB}{190, 0, 0}
\usepackage{gensymb}
